\documentclass[../../../include/open-logic-section]{subfiles}

\begin{document}

\olfileid{sth}{story}{rus}
\olsection{રસેલનો વિરોધાભાસ (ફરી)}

\olref[sfr][][]{part}માં આપણે સહજ ગણસિદ્ધાંત સાથે કામ કર્યું હતું.
પરંતુ \emph{અત્યંત} સહજ ખ્યાલ મુજબ ગણો માત્ર વિધેયધર્મોના વ્યાપો છે.
આ સહજ વિચાર નીચેનો સિદ્ધાંત ફરજિયાત બનાવે:

\begin{defish}
  \emph{સહજ ગણરચના.} કોઈ પણ સૂત્ર $\phi$ માટે $\Setabs{x}{\phi(x)}$ અસ્તિત્વ ધરાવે છે.
\end{defish}

આ સિદ્ધાંત આકર્ષક હોવા છતાં અસુસંગત છે, એ સાબિત કરી શકાય છે.
\olref[sfr][set][rus]{sec}માં આપણે આ જોયું હતું, પરંતુ આ પરિણામ
એટલું અગત્યનું અને એટલું સીધું છે કે તેને ફરી કહેવું યોગ્ય છે.
શબ્દશઃ.

\begin{thm}[રસેલનો વિરોધાભાસ]
$R = \Setabs{x}{x \notin x}$ એવો કોઈ ગણ નથી.
\end{thm}

\begin{proof}
જો $R = \Setabs{x}{x \notin x}$ અસ્તિત્વ ધરાવતો હોય, તો
$R \in R$ ત્યારે અને માત્ર ત્યારે જ્યારે $R \notin R$, જે પરસ્પરવિરોધ છે.
\end{proof}

રસેલે જૂન~1901માં આ પરિણામ શોધ્યું હતું. (જોકે તેમણે વિરોધાભાસને
અહીં રજૂ કરેલા સ્વરૂપમાં જ મૂક્યો નહોતો, કારણ કે તેઓ ફ્રેગેના
\emph{Grundgesetze}માં રૂપરેખા આપેલા ગણસિદ્ધાંતનો વિચાર કરતા હતા.
\olref[blv]{sec}માં આપણે આ પર પાછા આવીશું.) રસેલે 16~જૂન~1902ના
રોજ ફ્રેગેને પત્ર લખીને તેમની પદ્ધતિની અસુસંગતતા સમજાવી હતી.
આ પત્રવ્યવહાર અને થોડી પૃષ્ઠભૂમિ માટે
\citet[pp.~124--8]{Heijenoort1967} જુઓ.

આ બે પંક્તિની સાબિતી \emph{શુદ્ધ તર્કશાસ્ત્ર}નું પરિણામ છે એ પર
ભાર મૂકવો યોગ્ય છે. ખરું છે કે $\Setabs{x}{x \notin x}$ સંજ્ઞામાં
આપણે અપ્રત્યક્ષ રીતે ઘટકો દ્વારા નિર્ધારિત સમાનતાની (બિનતાર્કિક?)\
સ્વયંસિદ્ધિ વાપરી છે; કારણ કે $\Setabs{x}{\phi(x)}$ એ $\phi$
ધરાવતી વસ્તુઓનો \emph{અનન્ય} ગણ છે (આ સમાનતાના સિદ્ધાંતને કારણે),
જો એવો ગણ અસ્તિત્વ ધરાવતો હોય. પરંતુ પરિણામને નીચે પ્રમાણે કહીને
આ સિદ્ધાંતનો સંકેત પણ ટાળી શકીએ:
\emph{એવો કોઈ ગણ નથી જેના સભ્યો બરાબર સ્વસભ્ય ન હોય એવા ગણો જ હોય}.
અને આ વાતનો ગણો સાથે ખાસ સંબંધ નથી. રસેલે પોતે નોંધ્યું હતું તેમ,
બરાબર આ પ્રકારનું તર્ક આપણને નીચેના નિષ્કર્ષ પર લઈ જાય:
\emph{કોઈ પુરુષ બરાબર એવા પુરુષોની જ હજામત કરતો નથી જેઓ પોતાની
હજામત કરતા નથી}. અથવા: \emph{પગ જાતિનો કોઈ કૂતરો બરાબર એ જ જાતિના
એવા કૂતરાઓને જ સૂંઘતો નથી જેઓ પોતાને સૂંઘતા નથી}. અને એમ આગળ.
પ્રરૂપરૂપે પરિણામનું સ્વરૂપ માત્ર આ છે:
\[
\lnot \exists x \forall z(Rzx \liff \lnot R zz).
\]
અને આ તો પ્રથમ-ક્રમ તર્કશાસ્ત્રનું પ્રમેય (પ્રરૂપ) જ છે. તેથી
ગણસિદ્ધાંતમાં થોડોઘણો ફેરફાર કરીને રસેલનો વિરોધાભાસ ટાળી શકતા નથી;
આપણે ગણસિદ્ધાંત સુધી \emph{પહોંચીએ} તે પહેલાં જ તે ઊભો થાય છે.
જો આપણે (શાસ્ત્રીય) પ્રથમ-ક્રમ તર્કશાસ્ત્ર વાપરવાના હોઈએ, તો
$R = \Setabs{x}{x\notin x}$ એવો કોઈ ગણ નથી તે \emph{સ્વીકારવું}
જ પડે.

પરિણામ આ છે. જો અસુસંગતતા \emph{ટાળીને} સહજ ગણરચના સ્વીકારવી
હોય, તો માત્ર \emph{ગણસિદ્ધાંત}માં થોડોઘણો ફેરફાર કરવાથી કામ
ચાલશે નહીં. તેના બદલે \emph{તર્કશાસ્ત્ર}માં મૂળભૂત ફેરફાર કરવા પડે.

અલબત્ત, અશાસ્ત્રીય તર્કશાસ્ત્રો ધરાવતા ગણસિદ્ધાંતો રજૂ થયા છે.
પરંતુ ઓછામાં ઓછું એટલું તો કહી શકાય કે તેઓ માનક નથી. રસેલના
વિરોધાભાસ પ્રત્યેનો માનક અભિગમ તેને સીધી અનસ્તિત્વની સાબિતી તરીકે
લેવાનો અને પછી તેની સાથે કેવી રીતે કામ ચલાવવું તે શીખવાનો છે.
આપણે આ જ અભિગમ અપનાવીશું.

\end{document}
